\documentclass[11pt,letterpaper]{article}
\usepackage[margin=1in]{geometry}\usepackage{amsmath,amssymb,booktabs,array,microtype}\usepackage[T1]{fontenc}\usepackage{lmodern}
\usepackage[colorlinks=true,linkcolor=blue,citecolor=blue,urlcolor=blue]{hyperref}
\setlength{\parindent}{0pt}\setlength{\parskip}{5pt}\sloppy
\title{The Hubble-Kernel Growth Equation Against DESI DR1:\\ Consistent, Not Yet Discriminated, and the Precision That Decides}
\author{Carl P.\ Zimmerman\\ \small Briar Creek Tech\\ \small AI-assisted research programme; not peer reviewed}
\date{11 September 2026}
\begin{document}\maketitle
\begin{abstract}
The Hubble-kernel growth equation (doi:10.5281/zenodo.22706925) fixes the effective gravitational coupling of linear structure with no free parameter, $G_{\rm eff}(z)/G=\nu(cH(z)/a_0)$ with $(cH_0/a_0)^2=8\pi/(3\kappa^2\Omega_\Lambda)=49$, and predicts $f\sigma_8$ raised over Planck-$\Lambda$CDM by $2.7\%$, $1.9\%$, $1.1\%$ at $z=0.3,\,0.6,\,1.0$ (canonical footing; $3.8$, $2.6$, $1.6\%$ on the alternative footing). We confront it with the DESI DR1 full-shape growth measurements (DESI 2024 V, Table 9: ShapeFit $f\sigma_{s8}/(f\sigma_{s8})_{\rm fid}$ in six redshift bins, Planck fiducial, both the ShapeFit-only and ShapeFit+BAO sub-panels) and with the DESI modified-gravity parameter $\mu_0=0.05\pm0.22$. The equation is consistent with the data ($\chi^2$ per bin $0.65$--$0.73$ on both footings and both sub-panels) and is not yet discriminated from $\Lambda$CDM ($|\Delta\chi^2|<0.25$): DESI DR1's per-bin uncertainties of $10$--$19\%$ cannot resolve a $1$--$4\%$ effect. The equation's present-day coupling excess ($0.077$ canonical, $0.099$ alternative) lies $0.1$--$0.2\sigma$ from $\mu_0$. A $3\sigma$ test requires $\sim1\%$ on $f\sigma_8$ at $z\simeq0.3$--$0.6$, the DESI final-sample regime. This is a status result: the prediction is alive, on record before the decisive data, and its sign (more growth, not less) is the one the $S_8$ tension does not favour.
\end{abstract}
\section{The prediction}
With $a_0=\kappa c\sqrt{G\rho_\Lambda}$ and the Hubble-flow acceleration as the kernel's cosmological ambient field, $G_{\rm eff}(z)/G=\nu(cH(z)/a_0)$, $\nu(x)=1/(1-e^{-\sqrt x})$, $cH_0/a_0=7.0$ ($\kappa=\tfrac12$; $5.8$ on the alternative footing $a_0=1.13\times10^{-10}\,{\rm m\,s^{-2}}$). Linear growth with this coupling on a $\Lambda$CDM background ($\Omega_m=0.314$, $h=0.674$) gives the ratio of $f\sigma_8$ to its Planck-$\Lambda$CDM value listed in Table~\ref{t}.
\section{The data}
DESI 2024 V \cite{desiv} reports, for each of six redshift bins, the ShapeFit compressed parameter $f\sigma_{s8}/(f\sigma_{s8})_{\rm fid}$ (fiducial: Planck 2018 $\Lambda$CDM; MAP values with 68\% credible intervals), in two sub-panels: ShapeFit only, and ShapeFit with the post-reconstruction BAO information added. We symmetrise asymmetric intervals as the mean of the two sides. Separately, DESI 2024 \cite{desimg} constrains the phenomenological coupling $\mu(a)=1+\mu_0\,\Omega_\Lambda(a)/\Omega_\Lambda$ to $\mu_0=0.05\pm0.22$ from DESI full-shape and BAO combined with CMB, DES Y3 lensing and DES Y5 supernovae.
\section{Confrontation}
\begin{table}[h]\centering\caption{DESI DR1 ShapeFit growth ratios versus the equation.}\label{t}
\begin{tabular}{lcccccc}\toprule
bin & $z_{\rm eff}$ & ShapeFit only & ShapeFit+BAO & eq.\ (canonical) & eq.\ (alt) \\\midrule
BGS & 0.295 & $0.80\pm0.20$ & $0.84\pm0.19$ & 1.027 & 1.038\\
LRG1 & 0.510 & $1.09^{+0.12}_{-0.14}$ & $1.16\pm0.13$ & 1.021 & 1.029\\
LRG2 & 0.706 & $1.05\pm0.12$ & $1.04^{+0.11}_{-0.09}$ & 1.016 & 1.023\\
LRG3 & 0.919 & $0.96^{+0.11}_{-0.10}$ & $1.00^{+0.10}_{-0.08}$ & 1.012 & 1.018\\
ELG2 & 1.317 & $0.95^{+0.11}_{-0.08}$ & $0.95^{+0.10}_{-0.08}$ & 1.007 & 1.011\\
QSO & 1.491 & $1.16\pm0.12$ & $1.16\pm0.12$ & 1.006 & 1.009\\\bottomrule
\end{tabular}\end{table}
Over the six bins: $\chi^2=3.85$ ($\Lambda$CDM), $3.92$ (canonical), $3.98$ (alt) for ShapeFit only; $4.56$, $4.36$, $4.33$ for ShapeFit+BAO. The equation is consistent with the data ($\chi^2/6=0.65$--$0.73$) and undiscriminated from $\Lambda$CDM ($\Delta\chi^2$ between $-0.23$ and $+0.13$). Its present-day coupling excess, $G_{\rm eff}/G-1=0.077$ (canonical) or $0.099$ (alt), lies $0.12\sigma$ and $0.22\sigma$ from $\mu_0$.
\section{What decides}
The predicted effect is $1$--$4\%$ at $z<1$; DESI DR1 resolves $10$--$19\%$ per bin. A $3\sigma$ test needs $\sim1\%$ on $f\sigma_8$ at $z\simeq0.3$--$0.6$, which is the reach of the DESI final sample. A measured deficit of growth at the few-percent level, or a weak-lensing $S_8$ two to three percent below Planck-$\Lambda$CDM, falsifies the equation; an excess of the predicted size confirms it against $\Lambda$CDM. The prediction is on record before either. Provenance: repository lanes L180--L181 (\texttt{fable\_independent\_2026/}, committed scripts and outputs); the equation's derivation and scope are in doi:10.5281/zenodo.22706925.
\begin{thebibliography}{9}
\bibitem{desiv} DESI Collaboration, \emph{DESI 2024 V: Full-Shape Galaxy Clustering from Galaxies and Quasars}, arXiv:2411.12021.
\bibitem{desimg} M.~Ishak et al.\ (DESI Collaboration), \emph{Modified Gravity Constraints from the Full Shape Modeling of Clustering Measurements from DESI 2024}, arXiv:2411.12026.
\end{thebibliography}
\end{document}
